\section*{ВВЕДЕНИЕ}

\subsection*{Актуальность темы исследования}
\addcontentsline{toc}{subsection}{Актуальность темы исследования}

Современная защита от DDoS-атак, сетевая аналитика и телеметрия
операторских сетей всё чаще опираются на маршрутный контекст адресов.
Главный элемент этого контекста -- принадлежность IP-адреса автономной
системе: по ASN формируются правила фильтрации, выполняется скоринг
источников трафика, анализируются flow-записи и сопоставляются события с
операторской топологией.

Результаты работы внедряются в инфраструктуру продуктов BIFIT
MITIGATOR (защита от DDoS) и BIFIT COLLECTOR (анализ операторского
трафика). Оба продукта должны опираться на одну и ту же таблицу
\textit{IP-префикс~$\rightarrow$~ASN}, выведенную из полной таблицы
маршрутизации BGP; внешние геолокационные справочники с редким
обновлением не подходят для принятия решений в режиме реального времени.
Устаревшее соответствие префиксов и ASN может приводить к ошибкам
фильтрации, некорректной атрибуции источников атаки и расхождению
трактовки одного и того же события между продуктами. До начала работы
логика построения и обновления этой таблицы была встроена в BIFIT
COLLECTOR, а внутреннее хранилище \texttt{ranger} обновлялось по таймеру
полной заменой снимка: потоковое применение одиночных изменений к
in-memory базе не поддерживалось, а BIFIT MITIGATOR не мог
переиспользовать ту же реализацию без дублирования кода и данных.

Естественным первичным источником актуальных данных остаётся сам
протокол BGP~\cite{rfc4271}. Чтобы использовать его в сервисе
сопоставления префиксов и ASN, нужно получить начальный снимок RIB,
непрерывно обрабатывать
update-сообщения, поддерживать IPv4 и IPv6 и подобрать in-memory
структуру, которая одновременно обеспечивает малое время LPM-запроса и
позволяет применять инкрементальные изменения отдельных префиксных
записей без полной перестройки таблицы.

В режимах, где задержка актуализации на уровне нескольких часов не
влияет на прикладное решение, могут применяться внешние справочники.
Для высоконагруженных сервисов фильтрации и сетевой аналитики такой
лаг неприемлем: минуты расхождения локального состояния с глобальной
маршрутной таблицей повышают риск ложных срабатываний и пропусков
трафика~\cite{labovitz1997instability,labovitz2000convergence}.
Отсюда следуют требования к задержке актуализации на уровне единиц
секунд и к инкрементальному обновлению in-memory состояния вместо
периодической полной перестройки таблицы.

\subsection*{Степень разработанности темы}
\addcontentsline{toc}{subsection}{Степень разработанности темы}

В мировой литературе вопросы динамики BGP и стабильности маршрутов
изучаются с конца 1990-х. Работы Лабовица о маршрутной нестабильности и
времени сходимости BGP
\cite{labovitz1997instability,labovitz2000convergence} важны для данной
работы потому, что показывают: маршрутное состояние не является
статичным справочником и должно рассматриваться как изменяемый объект.
Формальная постановка задачи устойчивых путей у Гриффина и
др.~\cite{griffin2002stablepaths} задаёт теоретический язык для
обсуждения сходимости и переходных состояний. Эмпирические источники
RIPE~RIS, RouteViews и RIPE Atlas
\cite{ripe-ris,routeviews,ripe-beacons} релевантны как практическая база
для получения снимков и потоков маршрутных изменений.

Алгоритмическая сторона longest prefix match имеет столь же длинную
историю. PATRICIA-деревья~\cite{morrison1968patricia}, LC-trie
Нильссона и Карлссона~\cite{nilsson1999lctries}, controlled prefix
expansion Шринивасана~\cite{srinivasan1999prefixexpansion},
Tree~Bitmap Итертона~\cite{eatherton2004treebitmap} и Poptrie
Асаи~\cite{asai2015poptrie} рассматриваются как основные семейства
структур, с которыми необходимо сравнить выбранное хранилище. Их роль в
работе состоит не в прямом заимствовании реализации, а в задании
критериев сравнения: время LPM-запроса, стоимость обновления, расход
памяти, стоимость полной выгрузки и сложность сопровождения.
Adaptive~Radix~Tree предложен Лейсом и др.~\cite{leis2013art} для
in-memory индексов СУБД; в данной работе он рассматривается как
изменяемое хранилище префиксных ключей, а не как классический
маршрутизаторный алгоритм LPM.

Архитектурная парадигма «начальный снимок + поток обновлений» проработана
в литературе по системам обработки потоков
\cite{kreps2011kafka,akidau2015dataflow,carbone2015flink,kleppmann2017ddia}.
Эти источники вводятся для обоснования модели, в которой начальное
состояние строится из полного снимка, а последующие изменения применяются
как упорядоченная последовательность событий. В данной работе эта модель
адаптируется к BGP-производному состоянию и используется совместно с
выделением функции сопоставления «префикс~$\rightarrow$~ASN» в отдельный
сервис с общим контрактом для нескольких продуктов.

\subsection*{Решаемая проблема}
\addcontentsline{toc}{subsection}{Решаемая проблема}

Для продуктов DDoS-защиты и сетевой аналитики нужна единая точка истины
для пары \textit{префикс~$\rightarrow$~ASN}, доступная нескольким
потребителям и обновляемая с задержкой порядка единиц секунд. Прежняя
схема сочетала встроенность логики в один продукт, опрос по таймеру и
полную замену таблицы в \texttt{ranger}: при таком устройстве нельзя
получить ни общий сервис для BIFIT MITIGATOR и BIFIT COLLECTOR, ни
потоковое in-memory обновление без перестройки всей таблицы. Проблема,
решаемая в работе, состоит в выделении сопоставления в автономный
микросервис с общим gRPC API и в обеспечении штатного режима «начальный
снимок + поток обновлений» для in-memory состояния без отказа от
прикладно приемлемого времени чтения.

\subsection*{Объект и предмет исследования}
\addcontentsline{toc}{subsection}{Объект и предмет исследования}

Объектом исследования являются программные средства поддержки
актуального отображения сетевых префиксов в автономные системы по
данным BGP.

Предметом исследования являются архитектурные решения по выделению
функции сопоставления «префикс~$\rightarrow$~ASN» в отдельный
микросервис и методы потоковой актуализации его in-memory таблицы по
данным BGP, включая выбор изменяемого внутреннего хранилища,
поддерживающего LPM и инкрементальные правки.

\subsection*{Цель и задачи работы}
\addcontentsline{toc}{subsection}{Цель и задачи работы}

Цель работы -- разработать и исследовать микросервис, который
поддерживает актуальную таблицу соответствия
\textit{IP-префикс $\rightarrow$ ASN} по данным BGP для потребителей
DDoS-защиты и сетевой аналитики в режиме реального времени.

Для достижения цели в работе решаются такие задачи:

\begin{enumerate}
\item изучить предметную область междоменной маршрутизации, свойства
BGP и проблему свежести маршрутного состояния;
\item проанализировать известные подходы к LPM и структурам хранения
префиксных таблиц;
\item сформулировать функциональные и нефункциональные требования к
сервису, работающему по схеме «начальный снимок + поток обновлений»;
\item спроектировать архитектуру микросервиса и интерфейсы его
взаимодействия с источником маршрутных данных и потребителями;
\item реализовать загрузку начального снимка и потоковую обработку
обновлений с восстановлением после разрыва соединения;
\item реализовать шардированное in-memory хранилище на базе Adaptive
Radix Tree с операциями поиска, вставки, удаления и полной выгрузки;
\item провести сравнительный эксперимент с прежним вариантом и оценить
компромиссы между временем чтения, скоростью обновления и стоимостью
полной выгрузки.
\end{enumerate}

\subsection*{Методы исследования}
\addcontentsline{toc}{subsection}{Методы исследования}

В работе применяются методы системного анализа предметной области,
сравнительного анализа алгоритмов и структур данных, эксперимент в
форме воспроизводимых микробенчмарков на одном и том же стенде, а также
аналитическое моделирование совокупной CPU-нагрузки для двух режимов
работы. Каждый эксперимент по нескольку раз воспроизводится, а
полученные показатели сводятся в таблицы с одинаковым набором столбцов
для двух конкурирующих реализаций.

\subsection*{Методологическая и теоретическая основа исследования}
\addcontentsline{toc}{subsection}{Методологическая и теоретическая основа}

Методологическую базу составляют действующие RFC по BGP и его
расширениям~\cite{rfc4271,rfc4760,rfc4632,rfc2918,rfc7313,rfc7854}, а
также фундаментальные работы по динамике BGP
\cite{labovitz1997instability,labovitz2000convergence,
griffin2002stablepaths}. Теоретической основой алгоритмической части
служат работы по префиксным деревьям и LPM
\cite{morrison1968patricia,nilsson1999lctries,
waldvogel1997scalable,gupta1998hardware,srinivasan1999prefixexpansion,
eatherton2004treebitmap,asai2015poptrie}, а также оригинальная статья
по Adaptive Radix Tree~\cite{leis2013art}. Архитектурный шаблон
«начальный снимок + поток обновлений» взят из практики систем
обработки потоков и распределённых данных
\cite{kreps2011kafka,akidau2015dataflow,carbone2015flink,kleppmann2017ddia}.

\subsection*{Научная новизна}
\addcontentsline{toc}{subsection}{Научная новизна}

Научная новизна работы состоит не в новом сетевом протоколе, а в
постановке и решении прикладной задачи инфраструктурного уровня для
продуктовой линейки BIFIT. Во-первых, функция сопоставления
\textit{IP-префикс~$\rightarrow$~ASN} вынесена из состава BIFIT
COLLECTOR в автономный микросервис с единым gRPC контрактом для
нескольких потребителей -- это меняет границы ответственности и устраняет
необходимость дублировать логику между MITIGATOR и COLLECTOR. Во-вторых,
реализован переход от периодической полной замены таблицы к потоковому
применению BGP-событий к in-memory базе по шаблону «начальный снимок +
поток обновлений». В-третьих, для практической осуществимости потока
выбрано и внедрено шардированное in-memory ядро на базе
Adaptive~Radix~Tree~\cite{leis2013art}: структура обеспечивает
инкрементальное изменение отдельных префиксных записей, LPM-запросы,
полную выгрузку и рассылку подписчикам при эксплуатационно приемлемых
затратах CPU, хотя в маршрутной литературе ART чаще обсуждается в
контексте СУБД, а не BGP-состояния.

\subsection*{Положения, выносимые на защиту}
\addcontentsline{toc}{subsection}{Положения, выносимые на защиту}

\begin{enumerate}
\item Функция сопоставления «префикс $\rightarrow$ ASN» выделена в
отдельный микросервис с общим gRPC API для нескольких продуктов;
граница сервиса совпадает с границей предметной задачи, а не с границей
одного из потребителей.
\item Архитектурный шаблон «начальный снимок + поток обновлений»,
адаптированный к источнику BGP-данных, обеспечивает потоковое
поддержание in-memory таблицы без полной перестройки при каждом
изменении маршрута.
\item Шардированное in-memory хранилище на основе Adaptive~Radix~Tree
для пары «префикс $\rightarrow$ ASN» поддерживает LPM, точечные правки и
полную выгрузку и служит технической основой для пункта~2.
\item Алгоритм восстановления состояния после разрыва с источником,
сводящий пересинхронизацию к повторной загрузке снимка с последующим
переходом в режим потока.
\item Эмпирическая оценка преимущества потокового варианта по операции
инкрементального обновления и аналитическая модель совокупной
CPU-стоимости двух вариантов хранилища под профили нагрузки; численные
пределы, при которых пакетная схема перестаёт быть эксплуатационно
осуществимой.
\end{enumerate}

\subsection*{Степень достоверности и апробация результатов}
\addcontentsline{toc}{subsection}{Степень достоверности и апробация результатов}

Достоверность результатов обеспечивается тем, что все эксперименты
проведены на одном и том же стенде с фиксированной конфигурацией CPU,
оперативной памяти и версии среды исполнения Go, а измерения выполнены
стандартным механизмом \texttt{go test -bench} на нескольких размерах
таблицы и нескольких профилях нагрузки. Конфигурация стенда и набор
данных детально описаны в главе~\ref{sec:methodology}.
Апробация результатов проведена в инфраструктуре компании БИФИТ:
сервис включён в сборки продуктов BIFIT MITIGATOR и BIFIT COLLECTOR
и используется как единый источник пары «префикс $\rightarrow$ ASN»
для нескольких внутренних потребителей.

\subsection*{Практическая значимость}
\addcontentsline{toc}{subsection}{Практическая значимость}

Практическая значимость определяется применимостью результата как
инфраструктурного сервиса в системах сетевой аналитики, защиты
трафика и операторской отчётности, где важны:

\begin{itemize}
\item задержка актуализации таблицы маршрутов на уровне единиц секунд;
\item единая точка обслуживания нескольких клиентов через общий API;
\item согласованное и воспроизводимое состояние при перезапусках;
\item малая CPU-стоимость инкрементального обновления отдельных
префиксных записей;
\item унификация источника данных и снижение стоимости поддержки
нескольких продуктов компании.
\end{itemize}

\subsection*{Структура работы}
\addcontentsline{toc}{subsection}{Структура работы}

Работа содержит введение, десять глав основной части, заключение, список
использованных источников и список иллюстративного материала.
Главы~\ref{sec:bgp-domain}--\ref{sec:bgp-dynamics} посвящены
теоретическим основам BGP и динамике междоменной маршрутизации.
Глава~\ref{sec:storage-choice} содержит обзор структур LPM и обоснование
выбора in-memory хранилища префиксной таблицы. В
главе~\ref{sec:problem} формулируются задача исследования, гипотеза,
требования и критерии успеха. Главы~\ref{sec:architecture}--\ref{sec:art-storage}
описывают архитектуру микросервиса, алгоритмы потоковой актуализации и
реализацию шардированного ART-хранилища. Главы~\ref{sec:methodology} и
\ref{sec:experiment-results} посвящены методике и результатам
сравнительного эксперимента. В главе~\ref{sec:limitations} разобраны
ограничения текущей реализации и направления развития.

%==============================================================
